\documentclass[11pt,a4paper]{article}
\usepackage{helvet}
\renewcommand{\familydefault}{\sfdefault}
\usepackage{graphicx}
\usepackage{tikz}
\usepackage{eso-pic}
% Faded crest in the top-right corner of every page from where this is called
\newcommand{\cornercrest}{%
  \AddToShipoutPictureBG{%
    \begin{tikzpicture}[remember picture, overlay]
      \node[opacity=0.25, anchor=south east, xshift=-1cm, yshift=1cm]
        at (current page.south east)
        {\includegraphics[width=1.25cm]{figures/TCDnotext.png}};
    \end{tikzpicture}}}
\usepackage{amsmath}       % maths
\usepackage{parskip}
\usepackage[backend=biber, style=ieee]{biblatex}
\usepackage{hyperref}
\usepackage[top=3cm, bottom=3cm, left=3cm, right=3cm]{geometry}
\usepackage{setspace}
\setstretch{1.2}


\begin{document}

\begin{titlepage}
   \centering
   \vspace*{7cm}
  {\Huge\bfseries Formal Verification Group Project Report\par}
  {\Large Aidan McNally, Enda Benson, Evan Roe\par}
  \vspace{6cm}
  \includegraphics[width=0.35\textwidth]{figures/TCD.png}\par
  \vspace*{\fill}
\end{titlepage}

\cornercrest    

\section{Team Details}

- Names of team members
- Brief description of how you collaborated

Team members: Aidan McNally, Enda Benson, Evan Roe.

Collaborated through github repository. Each member attempted and contributed towards part 1, part 2 was started by Enda and completed by Aidan. Part 3 was completed by Evan and the report was a team effort.

\section{Use of AI Tools}

- Which tools (if any)
- How they contributed in this part of the project

No AI was used towards part 1, only the given resources for learning Dafny. Part 2 and consequently part 3 used AI to learn more about how Dafny evaluated methods and how it verified them, to better understand the use of invariants and lemmas to help Dafny. 

\section{Code Modifications}

- Original code snippet
- Modified code
- Whether the modification was suggested by an AI tool
- Rationale for the change (why it was needed / how it helped verification)
- Note: You do not need to track changes to invariants, assertions, ghost code, or other non-computational proof scaffolding.

\section{Reflection}

- Which methods were hardest and why
- How proof attempts changed your understanding of specification and correctness
- Whether AI helped or slowed you down
- Any other thoughts on the verification experience

For Part 3 the method that gave the most problems was the setProduct method. The nested loop and the 5 postconditions meant that Dafny needed to search through many loop iterations to verify the method. Many times we were faced with a verification timeout which meant we were not even informed about what part of the method could not be verified. The nested loop also meant that each loop needed its set of invariants, which did not come intuitively to duplicate the invariants over both loops.
AI definitely helped in this case, because since it was a new programming language essentially, learning its intricacies quickly was made simple with the use of AI. If we were stumped on a method in part 2 or 3, we could turn to AI to learn how Dafny evaluated and verified different postconditions types. Otherwise I expect we would have needed to dig through Dafny documentation to get there.

This report does not need to be long, but should be of reasonable quality.

You should submit to blackboard two files before the deadline Dec. 01, 2026
Even-Odd-Sets-part3.dfy 
2/3 of Part 3 mark distributed on successful verification of each given pre-/post-condition
Your report in PDF format
1/3 of Part 3 mark

\end{document}